\documentclass[11pt]{article}
\usepackage{lmodern}
\usepackage[T1]{fontenc}
\usepackage[letterpaper,margin=0.8in]{geometry}
\usepackage{xcolor}
\usepackage{booktabs}
\usepackage{tabularx}
\usepackage{enumitem}
\usepackage{array}
\usepackage[hidelinks]{hyperref}

\setlist{itemsep=1pt,topsep=3pt}
\newcommand{\br}[1]{\texttt{#1}}
\newcommand{\h}[1]{\texttt{#1}}
\newcolumntype{L}{>{\raggedright\arraybackslash}X}

\title{PACMAN Zipline Firmware and Software (3.X)\\[2pt]
  \large Branches, builds, boards and status}
\author{Maintained by the PACMAN agents (augie)}
\date{2026-11-04}

\begin{document}
\maketitle
\sloppy

\section{Purpose and scope}

Zipline (3.X) is the current PACMAN firmware and software, for the 1v5 board (also built for the 1v4), after the
legacy (1.X) and hijinks (2.X) generations (companion notes \h{legacy.tex}, \h{hijinks.tex}).  This note records
where zipline stands: the branches, what was built from which commit, which build is on which board, the results,
and what is open.  The firmware itself (registers, units, timing) is described in \h{firmware.tex}; the checkout
tools in \h{qaqc.tex}.  Digests are SHA-256, first eight hex digits.

\section{Repositories}

\begin{center}
\begin{tabularx}{\linewidth}{@{}llL@{}}
\toprule
Repository & Branch & State \\
\midrule
\h{mulhearn/pacman-firmware} & \br{zipline-2023.2} & tag \br{v3.6} (a85681c, the release); branch at aaa46c1
  (bump to 3.7, candidate), also \br{working-on-3.7-zipline-2023.2} \\
\h{mulhearn/pacman-software} & \br{zipline-2023.2} & tags \br{v3.6.0} (94e3f01) and \br{v3.6.1} (90d4e19); branch
  at 90677a7 (bump to 3.6.2, candidate), also \br{working-on-v3.6.2-zipline-2023.2} \\
\bottomrule
\end{tabularx}
\end{center}

\begin{itemize}
\item Toolchain: Vivado 2023.2 (firmware), PetaLinux 2023.2 (images and the Linux applications), Vitis 2023.2 (bare
  metal).
\item Firmware simulation: GHDL, \h{ghdl/} (test benches) and \h{ghdl/agents} (\h{make check} against
  \h{docs/requirements.txt}); \h{vivado/src/hdl} is kept in step with \h{ghdl/} in each commit.
\item Software: the Linux applications in \h{petalinux/pkg/pacman-apps} (\h{pacman\_cmdserver},
  \h{pacman\_dataserver}, \h{pacman\_menu}, \h{pacman\_hwcheck}, \h{pacman\_fwcheck}, \h{pacman\_lnxcheck}), one
  copy of the drivers (\h{drivers-common}) shared with the bare-metal menu (\h{apps/menu}).
\end{itemize}

\section{Firmware 3.6}

3.6 adds to 3.5: the ADC unit, the logic analyzer with ADC capture and cross triggers, the \h{\_UPPER}/\h{\_LOWER}
register names, the RX packet timeout in units of 256 cycles, and the RX sweep that resumes at the next channel after
a packet closed by its maximum.  Released 2026-11-02 as 3.6 (a85681c, \h{BUILD\_RELEASE}; tag \br{v3.6}); the
next commit bumps the version to 3.7 (candidate).

\begin{center}
\begin{tabularx}{\linewidth}{@{}lllL@{}}
\toprule
Commit & Date & Build & Result \\
\midrule
650dea6 & 2026-10-26 & hw1v5 & bit 82c2be63, BOOT.BIN c43b1340; timing met (WNS +0.115\,ns, WHS +0.049\,ns);
  \h{pacman\_fwcheck} PASS (30 of 30 requirements, twice) \\
ae71967 & 2026-10-29 & --- & on hw1v4, a reset on tiles 1--8 asserts the single legacy reset (was 1--10) \\
928fcf0 & 2026-10-30 & hw1v4 & bit 2755aa4c; timing met (WNS +0.185\,ns, WHS +0.011\,ns) \\
a85681c & 2026-11-02 & hw1v5 & the release, from a fresh clone: bit 6d84be5d, official (DIRTY 0); timing met
  (WNS +0.117\,ns, WHS +0.014\,ns), no critical warnings \\
a85681c & 2026-11-02 & hw1v4 & the release: bit aa06b896, official; timing met (WNS +0.075\,ns, WHS +0.025\,ns);
  one critical warning (ADC inputs unconnected in the hw1v4 block design, open) \\
\bottomrule
\end{tabularx}
\end{center}

\noindent The release bitstreams are in the release images.  Utilization of the release (hw1v5, post-place,
xc7z020): 30\,725 slice LUTs (57.8\%), 50\,712 registers (47.7\%), 45 block RAM tiles (32.1\%); hijinks 2.8 used
39\,788 LUTs (74.8\%), 60\,644 registers (57.0\%) and 99 tiles (70.7\%), the v1.5 1.X line (\br{sync-fix}) 40\,777
LUTs (76.7\%).  The release still holds a VIO (109 LUTs) and the Vivado debug hub it brings in (449 LUTs).

\section{Releases}

\begin{itemize}
\item A release is a dedicated commit that sets only the release flag (firmware \h{BUILD\_RELEASE},
  software \h{PACMAN\_SOFTWARE\_RELEASE}~1); the next commit bumps the version and clears the flag, nothing else.
  Builds of the release commit report ``official'', all others ``candidate''.
\item Software 3.6.0 (94e3f01, 2026-11-03): Michael's production build, with the official site overlay, on the 1v5:
  \h{pacman\_lnxcheck} 23 of 23, \h{pacman\_fwcheck} 30 of 30, \h{pacman\_hwcheck} all PASS, loopback 0 errors in
  5.1\,Mbit; LArPix v3 configured by larpix-lib through the servers (\h{larpix-lib.tex}).
\item Software 3.6.1 (90d4e19, 2026-11-03): the ASIC power scripts (\h{power\_up\_tile.sh}, \h{power\_down.sh},
  \h{report\_power.sh}) ported to zipline, through the command server at the \h{pacman\_menu} nominal values;
  package scripts for the \h{pacman-apps} tarball (software-only updates, untarred at \h{/}) and the
  \h{firmware-dev} kit (the \h{BOOT.BIN} inputs and a \h{bootgen.bif}, for firmware developers).
\item Point releases confined to \h{pacman-apps} ship only its tarball, no new images.  The bundle
  (Michael's \h{bundle-v3.6.0}) is named for the firmware: the 3.6.0 images (\h{BOOT.BIN} 801a5711,
  \h{boot.scr} 3b9b939d, \h{image.ub} 8094cb12), the 3.6.1 root filesystem 48e50939, \h{pacman-apps-v3\_6\_1}
  2a67b0a4, \h{firmware-dev-v3\_6\_1} 8c0ab6b1.  The 3.6.1 images match 3.6.0 bit for bit except the root filesystem.
\end{itemize}

\section{Software and images}

Recent changes (pacman-software):
\begin{itemize}
\item 3d77617 to f516bc7 (2026-10-30): support for the ADC and logic analyzer; \h{pacman\_hwcheck} as a test suite
  and \h{pacman\_fwcheck}; the I2C overhaul (sticky status, bus recovery, MUX refusals as returns); one copy of the
  drivers; the I2C stress test.
\item 2026-10-31: I2C at 100\,kHz by default (640a124, needed on a hw1v4 card); ASIC TX at full drive in the
  configuration root (6d2ad8e); \h{-{}-sma} tolerance 20\% (41c6e00); the I2C device skip removed (3c9b2d0).
\item 2026-11-01: the \h{pacman\_hwcheck} tests as a submenu of both main menus (da113dc); one menu input on both
  platforms, a letter then Enter in every menu, a clean end at the end of input; one copy of the RX/TX and ADC menus for both platforms, the power menu folded
  into the I2C menu, a 16-BD ring for the menus (671a3dd);  on the 1v4 since 2026-11-01.
\item 2026-11-01 (290979d, which replaces 671a3dd with the changes below added): one RX/TX loopback core, which uses the configuration and ring as it finds them; in the
  RX/TX menu the ring size (16 to 512 BDs, default 512, set up by \h{m}), the loopback length, the RXTX UART mask
  (including tiles 2+3) and a BER check; \h{pacman\_hwcheck -{}-loopback} per tile, then all the masked tiles at
  once; the RX config and RX buffer config words built by functions.  On the 1v4 since 2026-11-01.
\item 2026-11-01 evening (4e9164c, 9a36407, 9e30354): \h{pacman\_lnxcheck} memory limits 700--800\,MB (1\,GB less
  the DMA buffer); menu presets for the bench cards on tiles 6+7 and 9; the ramdisk and device tree loaded below the
  DMA buffer (U-Boot placed the ramdisk inside it, so the buffer's reservation failed); no PXE tries in the TFTP
  fall-through, and \h{boot.scr.uimg} copied into \h{images/linux}; the Python library on Python's path in every
  context (\h{pacman.pth}; \h{pacman\_setup.sh} removed); root may log in over ssh with a password, once one is set;
  a site overlay: the untracked \h{petalinux/site} (files copied over the root filesystem, then an optional
  \h{post.sh} for edits such as the root password hash) applied by the in-place builds, with \h{petalinux/site.example}
  as the template.  Tested on the 1v5.
\end{itemize}

\noindent The PetaLinux projects (\h{petalinux/ramdisk}, \h{petalinux/full}) are built in place, incrementally,
after each software or firmware change:

\begin{center}
\begin{tabular}{@{}lllll@{}}
\toprule
Software & Date & Image & \h{BOOT.BIN} & \h{image.ub} \\
\midrule
f516bc7 & 2026-10-30 & ramdisk & 41851385 & 7d452b37 \\
3c9b2d0 & 2026-10-31 & ramdisk & 12813299 & d7761858 \\
3c9b2d0 & 2026-10-31 & full & 4c1d3a0e & 8094cb12 \\
671a3dd & 2026-11-01 & ramdisk & cec3182c & c4bd61b8 \\
671a3dd & 2026-11-01 & full & 96d3faca & 8094cb12 \\
290979d & 2026-11-01 & ramdisk & 5bf59506 & 73368917 \\
290979d & 2026-11-01 & full & 348baf8a & 8094cb12 \\
9e30354 & 2026-11-01 & ramdisk & 31e01844 & 7b229452 \\
9e30354 & 2026-11-01 & full & 8998d25c & 8094cb12 \\
\bottomrule
\end{tabular}
\end{center}

\section{Boards}

\begin{center}
\begin{tabularx}{\linewidth}{@{}lllL@{}}
\toprule
Board & Firmware & Image & Results \\
\midrule
1v5 (PACMAN 51 UCD) & a85681c & 3.6.0 (SD) & Full SD image of the 3.6.0 release (root on p2), \h{pacman-apps}
  3.6.1 installed by tarball.  2026-11-03, 3.6.1: ASIC configuration via \h{pacman\_menu}; tile power and monitoring via
  \h{pacman\_menu} and the power scripts (tile 1: VDDA 2365\,mV, VDDD 1301\,mV, matching the menu); SYNC words;
  \h{pacman\_units}; RX/TX internal loopback from \h{pacman\_cmdserver} to \h{pacman\_dataserver}
  (\h{pacman\_push}, \h{pacman\_counter}, e.g.\ \h{-w 400 -n 100 -d 0}); RX test patterns at full speed on all 40
  UARTs, 6060 packets/ms sustained (the theoretical maximum: 100\,MHz / (10 clocks $\times$ 66 bits) $\times$ 40) with
  no loss, \h{pacman\_dataserver} at 70\% CPU.  Earlier (650dea6, 290979d): 2026-10-30, f516bc7:
  \h{pacman\_hwcheck}, every test PASS on the tiles carrying the test cards, except \h{-{}-ana\_mon} on tile 4 (a weak
  TESTP leg, set aside).  2026-11-01, 290979d: \h{pacman\_lnxcheck} 23 of 23 PASS (memory limits 700--800\,MB,
  4e9164c), \h{pacman\_fwcheck} 30 of 30 PASS, \h{pacman\_hwcheck} PASS on tiles 6+7 (loopback card), 9 (ATC card)
  and the board-wide tests; the production SD boot with TFTP fall-through PASS; LArPix v3 on tile 1 replies and
  passes the ASIC loopback (\h{larpix.tex}).  Bare-metal menu (4e9164c): boot, versions, I2C, hwcheck read-only,
  ADC, I2C stress and the ASIC hello PASS; the RX/TX and ASIC loopbacks lose words (open) \\
1v4 & 928fcf0 & 290979d (TFTP) & 2026-11-01: the new menus work (letters and Enter, a clean end at the end of
  input, the hwcheck submenu: read only as expected, I2C stress PASS); the LArPix v3 on tile 1 does not reply
  (UART 0 streams junk; it replied exactly on 2026-10-31 with 3c9b2d0): the configuration sent right after the full
  reset is lost, and a 0.2\,s pause fixes it.  2026-11-01, 290979d: a 16-BD ring set up from the menu, the ASIC hello
  exact twice, and the BER loopback on tiles 2+3 (internal loopback): 80000 of 80000 words, 0 bit errors in
  5.12\,Mbit \\
\bottomrule
\end{tabularx}
\end{center}

\noindent The 1v4 has no SD card: it boots by JTAG and TFTP; it is off the bench since 2026-11-01.  On the 1v5, ssh
and scp to the board work by key (sshd allows root by key only).

\section{Open}

\begin{itemize}
\item ZMQ (queued for 3.7): whole messages can be dropped between \h{pacman\_dataserver} and a client without
  any error (seen with \h{pacman\_push -w 40} at a long delay; \h{-w 400 -d 0} counts all).  Suspect: the
  fire-and-forget PUB socket (high-water mark 1, send timeout 0), on which \h{zmq\_msg\_send} succeeds although the
  message is dropped.  \h{pacman\_counter} prints no final report and not its message count.  In the 3.6.1 release
  notes as a known issue.
\item The VIO in the release (in discussion for 3.7).
\item \h{pacman\_menu} test patterns: the payload is written after the enable, so the first frame carries the old
  payload; \h{0x89ABDCDF} is probably meant as \h{0x89ABCDEF}.
\item Bare metal: a small DMA ring (16 BDs) loses words in the loopback tests (Linux copes); the default 512 is fine.
\item Planned for 3.7: registers for communication between processes (an I2C error flag, stalling the dataserver,
  lost-packet reports, the DMA race), polling for the end of a reset, more robust I2C (a termination option, a richer
  simulation), the \h{pattern.vhd} start-up.
\item Robustness of the Vivado build: \h{create.tcl} adds the HDL from the repository, so HDL edited after the
  project was created partly enters a rebuild; release bitstreams are built from a fresh clone.
\end{itemize}

\end{document}
